% BHU PYQ -- Manually transcribed & error-corrected
% Paper: BCF-214 : Corporate Auditing
% Exam: B.Com. (Hons.) FMM Semester III Examination, 2013-14

\documentclass[11pt,a4paper]{article}
\usepackage[a4paper,top=2.5cm,bottom=2.5cm,left=2.8cm,right=2.8cm,headheight=14pt]{geometry}
\usepackage{amsmath,amssymb}
\usepackage[T1]{fontenc}
\usepackage[utf8]{inputenc}
\usepackage{lmodern,microtype,fancyhdr,enumitem,hyperref,lastpage,parskip}

\hypersetup{
    colorlinks=true,
    urlcolor=black,
    linkcolor=black,
    pdftitle={BCF-214: Corporate Auditing -- BHU B.Com. (Hons.) FMM Sem III 2013-14}
}

\pagestyle{fancy}
\fancyhf{}
\renewcommand{\headrulewidth}{0.4pt}
\renewcommand{\footrulewidth}{0.4pt}
\fancyhead[L]{\small   \textit{Banaras Hindu University}}
\fancyhead[R]{\small   \textit{BCF-214: Corporate Auditing}}
\fancyfoot[C]{\small Page  \thepage\ of \pageref{LastPage}}


\newlist{parts}{enumerate}{2}
\setlist[parts,1]{label=(\alph*),leftmargin=*,itemsep=5pt,topsep=3pt}
\setlist[parts,2]{label=(\roman*),leftmargin=*,itemsep=3pt,topsep=2pt}

\newcommand{\pts}[1]{\hfill   \textit{[#1]}}


\begin{document}


\begin{center}
  {\large\bfseries BANARAS HINDU UNIVERSITY}\\[2pt]
  {\normalsize B.Com. (Hons.) FMM --- Semester III Examination, 2013-14}\\[6pt]
  \rule{0.8\linewidth}{0.5pt}\\[6pt]
  {\large\bfseries Paper No. BCF-214: Corporate Auditing}\\[4pt]
  {\normalsize Time: 3 Hours\hfill Maximum Marks: 70}
\end{center}

\rule{\linewidth}{0.4pt}

\smallskip



\noindent   \textit{Instructions: Answer all questions. All questions carry equal marks (14 marks each).}

\bigskip




\noindent   \textbf{Question 1.} \\
Explain the different types of errors and frauds found in auditing a firm's accounts. Can an auditor prevent such errors and frauds?\pts{14}

\centerline{   \textbf{OR}}

Define continuous audit. Discuss its merits and demerits. For what type of business it is suitable?\pts{14}

\medskip\rule{\linewidth}{0.2pt}




\noindent   \textbf{Question 2.} \\
Before conducting a new audit, discuss the steps to be taken by an auditor.\pts{14}

\centerline{   \textbf{OR}}

What is Audit Notebook? Describe its contents. How far it is useful to an auditor?\pts{14}

\medskip\rule{\linewidth}{0.2pt}




\noindent   \textbf{Question 3.} \\
What points should be kept in mind while examining the vouchers and vouching payments?\pts{14}

\centerline{   \textbf{OR}}

How would you vouch the following items?\pts{7+7}

\begin{parts}
  \item Purchase of machinery
  \item Cash sales
\end{parts}

\medskip\rule{\linewidth}{0.2pt}




\noindent   \textbf{Question 4.} \\
What do you mean by valuation of assets? Explain its objects.\pts{14}

\centerline{   \textbf{OR}}

How would you verify the following?\pts{7+7}

\begin{parts}
  \item Investment
  \item Plant and machinery
\end{parts}

\medskip\rule{\linewidth}{0.2pt}




\noindent   \textbf{Question 5.} \\
Discuss the provisions of Companies Act regarding appointment of company auditor.\pts{14}

\centerline{   \textbf{OR}}

Explain the civil liabilities of a company auditor.\pts{14}

\vfill
\rule{\linewidth}{0.4pt}

\begin{center}\small
\href{https://bhu-exam.vercel.app/}{Click here to visit our website}\\[4pt]
\href{https://me-aryan.vercel.app/}{Click here to visit our contributor}\\[4pt]
\href{https://quashan-me.vercel.app/}{Click here to visit our contributor}
\end{center}

\end{document}
